\documentclass[12pt]{article}
\usepackage[utf8]{inputenc}
\usepackage[T1]{fontenc}
\usepackage[margin=1in]{geometry}
\usepackage{setspace}
\usepackage{xcolor}
\usepackage[colorlinks=true, urlcolor=blue!60!black, linkcolor=blue!60!black]{hyperref}
\setlength{\parindent}{2em}
\setlength{\parskip}{1.25em}
\renewcommand{\baselinestretch}{1.0}
\begin{document}
\begin{center}
{\LARGE \textit{United States v. Morrison}, 529 U.S. 598 (2000)\par}
\vspace{0.8em}
{\large Case Brief --- Four-Page Edition\par}
\end{center}
\vspace{1em}

\section{Memory Jogger}

Applying \textit{Lopez} and the state-action requirement, the Court struck down the civil remedy of the Violence Against Women Act: Congress may not regulate noneconomic violent crime based on aggregate economic effects, even with voluminous findings, nor remedy state bias with a nationwide cause of action against private individuals (\href{https://www.courtlistener.com/opinion/118363/united-states-v-morrison/}{CourtListener}).

\section{Detailed Case Facts}

Christy Brzonkala, a Virginia Tech student, alleged that within 30 minutes of meeting them in September 1994, football players Antonio Morrison and James Crawford assaulted and repeatedly raped her, and that Morrison later boasted about it in crude terms. She became depressed, stopped attending classes, and withdrew. At a university hearing, Morrison admitted sexual contact despite her twice saying ``no.'' He was suspended for two semesters, but after a second hearing the offense was recast as ``using abusive language,'' and the provost set aside the suspension as ``excessive.'' Learning from a newspaper that Morrison would return, Brzonkala dropped out.

VAWA's 42 U.S.C. \S~13981 declared a right to be free from ``crimes of violence motivated by gender'' and created a federal civil action for damages and injunctive relief against perpetrators. Congress invoked the Commerce Clause and Section 5, citing four years of hearings and task-force reports from 21 States documenting the economic toll of such violence and gender bias in state justice systems.

\section{Procedural History}

Brzonkala sued Morrison, Crawford, and Virginia Tech in the Western District of Virginia; the United States intervened to defend \S~13981. The district court held that Congress lacked power to enact it. A Fourth Circuit panel reversed, but the en banc court affirmed the district court. The Supreme Court granted certiorari, heard argument on January 11, 2000, and affirmed on May 15, 2000.

\section{Judicial Votes}

Affirmed, \textbf{5--4}.
\begin{itemize}
\item \textbf{Majority:} Chief Justice Rehnquist, joined by Justices O'Connor, Scalia, Kennedy, and Thomas.
\item \textbf{Concurrence:} Justice Thomas.
\item \textbf{Dissents:} Justice Souter, joined by Justices Stevens, Ginsburg, and Breyer; Justice Breyer, joined by Justice Stevens (and by Souter and Ginsburg as to Part I-A).
\end{itemize}

\section{Holding}

Section 13981 exceeds both the Commerce Clause and Section 5. Congress may not regulate noneconomic, violent conduct based solely on its aggregate effect on interstate commerce, even with extensive findings resting on a but-for causal chain. Section 5 reaches only state action and requires ``congruence and proportionality''; a nationwide remedy against private individuals exceeds it. The Court left intact VAWA's interstate-travel crime and funding programs and declined to adopt a categorical rule against aggregating noneconomic activity.

\section{Analysis of Opinions}

\textbf{Majority (Rehnquist).} Applying the \textit{Lopez} factors, Rehnquist held that gender-motivated violence is ``not, in any sense of the phrase, economic activity''; that \S~13981 had no jurisdictional element; and that Congress's findings were ``substantially weakened'' because they relied on the but-for reasoning \textit{Lopez} rejected, which would let Congress regulate murder, marriage, and childrearing. On Section 5, the \textit{Civil Rights Cases} and \textit{United States v. Harris} confine the Amendment to state action; \textit{Guest} and \textit{Carter} were dicta. The remedy visited no consequence on any state official and applied nationwide though bias was not shown in most States, failing \textit{City of Boerne}.

\textbf{Thomas (concurring).} The ``substantial effects'' test itself is ``rootless and malleable'' and should be replaced with an originalist standard.

\textbf{Souter (dissenting).} The decision cannot be squared with \textit{Wickard} or \textit{Heart of Atlanta}. Congress's record, showing that domestic violence costs \$3 to \$10 billion a year and that nearly half of rape victims lose their jobs, satisfies rational basis; the majority revives pre-1937 formalism. With 36 States supporting the law, Morrison had ``won the states' rights plea against the states themselves.''

\textbf{Breyer (dissenting).} The economic/noneconomic line is unworkable (is a mugger engaged in economic activity?), Congress must strike the federal balance, and a federal substitute remedy for documented state-court bias is a proportionate Section 5 measure.

\section{Examples: Future Applications}

\begin{enumerate}
\item \textbf{Same-Side (Invalid): Hate-Crime Tort.} A federal damages action for any bias-motivated assault, resting on aggregate findings with no jurisdictional element, fails under both powers.
\item \textbf{Same-Side (Invalid): Arson of a Private Home.} Federal arson charges for burning an owner-occupied, noncommercial home cannot rest on aggregate insurance reasoning (compare \textit{Jones v. United States}).
\item \textbf{Opposite-Side (Valid): Interstate Domestic Violence.} Crossing state lines to injure an intimate partner, 18 U.S.C. \S~2261(a)(1), regulates the channels of commerce, as footnote 5 noted.
\item \textbf{Opposite-Side (Valid): Suits Against State Police.} A statute targeting departments with documented equal-protection violations reaches state action and can be congruent and proportional.
\item \textbf{Fence-Sitter: Intrastate Endangered Species.} Barring brush-clearing that destroys a noncommercial species' habitat looks local, yet may be part of a larger economic scheme under \textit{Gonzales v. Raich}.
\end{enumerate}

\section{Critique}

\textbf{Originalist views.} Randy Barnett praised \textit{Morrison} for showing \textit{Lopez} was no fluke and for defending enumerated powers against limitless aggregation. Strict originalists, like Thomas, faulted the Court for keeping the substantial-effects test, and some historians question reliance on the 1883 \textit{Civil Rights Cases} as faithful to Reconstruction.

\textbf{Progressive and feminist views.} Robert Post and Reva Siegel argued that \textit{Morrison} paired a constricted Commerce Clause with a judicial monopoly over Fourteenth Amendment enforcement, treating gender violence as ``local.'' Kermit Roosevelt argued the Section 5 holding misapplied state action to a remedy aimed at state-court failures, and Erwin Chemerinsky questioned whether judicial line-drawing protects liberty.

\section{Key Quotations}

\begin{quote}
``We accordingly reject the argument that Congress may regulate noneconomic, violent criminal conduct based solely on that conduct's aggregate effect on interstate commerce. The Constitution requires a distinction between what is truly national and what is truly local.''
\end{quote}
\noindent --- Chief Justice Rehnquist, Opinion of the Court, 529 U.S. 598.

\begin{quote}
``Until this Court replaces its existing Commerce Clause jurisprudence with a standard more consistent with the original understanding, we will continue to see Congress appropriating state police powers under the guise of regulating commerce.''
\end{quote}
\noindent --- Justice Thomas, Concurring, 529 U.S. 598.

\begin{quote}
``One obvious difference from United States v. Lopez, 514 U.S. 549 (1995), is the mountain of data assembled by Congress, here showing the effects of violence against women on interstate commerce.''
\end{quote}
\noindent --- Justice Souter, Dissenting, 529 U.S. 598.

\end{document}
